\vspsec
\section{Related Work}
\vspsec
Advances in learning control policies have shown success on numerous complex and high dimensional tasks~\citep{schulman2015trust,lillicrap2015ddpg,mnih2015human,levine2016end,silver2017mastering}. While reinforcement learning algorithms provide a framework for learning new tasks, they primarily focus on mastery of individual skills, rather than generalizing and quickly adapting to new scenarios.
Furthermore, model-free approaches~\citep{peters2008reinforcement} require large amounts of system interaction to learn successful control policies, which often makes them impractical for real-world systems. In contrast, model-based methods attain superior sample efficiency by first learning a model of system dynamics, and then using that model to optimize a policy~\citep{deisenroth2013survey,lenz2015deepmpc,levine2016end,nagabandi2017roach,williams2017information}.  Our approach alleviates the need to learn a single global model by allowing the model to be adapted automatically to different scenarios online based on recent observations. A key challenge with model-based RL approaches is the difficulty of learning a global model that is accurate for the entire state space. Prior model-based approaches tackled this problem by incorporating model uncertainty using Gaussian Processes (GPs)~\citep{ko2009gp,deisenroth2011pilco, doerr2017mbpid}. However, these methods make additional assumptions on the system (such as smoothness), and does not scale to high dimensional environments.  \citet{chua2018deep} has recently showed that neural networks models can also benefit from incorporating uncertainty, and it can lead to model-based methods that attain model-free performance with a significant reduction on sample complexity. Our approach is orthogonal to theirs, and can benefit from incorporating such uncertainty.

Prior online adaptation approaches~\citep{tanaskovic2013adaptive,aswani2012extensions} have aimed to learn an approximate global model and then adapt it at test time. Dynamic evaluation algorithms~\citep{rei2015onlinerepre,krause2017dynamiceval,krause2016multiplicative,fortunato2017bayesian}, for example, learn an approximate global distribution at training time and adapt those model parameters at test time to fit the current local distribution via gradient descent. There exists extensive prior work on online adaptation in model-based reinforcement learning and adaptive control~\citep{sastry1989adaptive}. In contrast from inverse model adaptation~\citep{kelouwani2012online,underwood2010online, pastor2011online, franzi2016robustonline, franzi2016drift, rai2017feedback}, we are concerned in the problem of adapting the forward model, closely related to online system identification~\citep{manganiello2014optimization}. Work in model adaptation~\citep{levine2013guided,gu2016continuous,fu2015onlineadapt, weinstein2017motorcontrol} has shown that a perfect global model is not necessary, and prior knowledge can be fine-tuned to handle small changes. These methods, however, face a mismatch between what the model is trained for and how it is used at test time. In this paper, we bridge this gap by explicitly training a model for fast and effective adaptation. As a result, our model achieves more effective adaptation compared to these prior works, as validated in our experiments.

Our problem setting relates to meta-learning, a long-standing problem of interest in machine learning that is concerned with enabling artificial agents to efficiently learn new tasks by learning to learn~\citep{thrun1998learning,schmidhuber1991learning,naik1992meta,lake2015human}. A meta-learner can control learning through approaches such as deciding the learner's architecture~\citep{baker2016designing}, or by prescribing an optimization algorithm or update rule for the learner~\citep{bengio1990learning,schmidhuber1992learning,younger2001meta,andrychowicz2016learntolearn,li2016learning,ravi2016optimization}.
Another popular meta-learning approach involves simply unrolling a recurrent neural network (RNN) that ingests the data~\citep{santoro2016one,munkhdalai2017meta,munkhdalai2017learning,mishra2017simple} and learns internal representations of the algorithms themselves, one instantiation of our approach (ReBAL) builds on top of these methods.
On the other hand, the other instantiation of our method (GrBAL) builds on top of MAML~\citep{finn2017maml}. GrBAL differs from the supervised version of MAML in that MAML assumes access to a hand-designed distribution of tasks. Instead, one of our primary contributions is the online formulation of meta-learning, where tasks correspond to temporal segments, enabling ``tasks'' to be constructed automatically from the experience in the environment.

Meta-learning in the context of reinforcement learning has largely focused on model-free approaches~\citep{duan2016rl2,wang2016learning,sung2017learning,al2017cont}. However, these algorithms present even more (meta-)training sample complexity than non-meta model-free RL methods, which precludes them from real-world applications. Recent work~\citep{saemundsson2018meta}  has developed a model-based meta RL algorithm, framing meta-learning as a hierarchical latent variable model, training for episodic adaptation to dynamics changes; the modeling is done with GPs, and results are shown on the cart-pole and double-pendulum agents. In contrast, we propose an approach for learning online adaptation of high-capacity neural network dynamics models; we present two instantiations of this general approach and show results on both simulated agents and a real legged robot.